\addtocontents{toc}{\protect\Needspace{3\baselineskip}}
\originalchapter{三次测验的阅读题与旋转坐标题}
本章补齐 8.286 Fall 2013 三次正式测验此前未收录的四个题组。题面按公开测验和官方解答改写；涉及 Weinberg 与 Ryden 书中插图时，用文字交代解题所需的图意，不复制受限书页。数值答案属于 2013 年课程的历史口径，不视作新的参数测量。

\src{8.286 Fall 2013 Quiz 1, Q1(a)--(e), 25分；官方解答逐项可得。}
\begin{exercise}
(a) Weinberg 用以 A、B、C 星系为参照的 Hubble 流示意图说明哪一项宇宙性质？解释为何从 A 看四周退行并不使 A 成为特殊中心。
(b) 约 $3\,\mathrm K$ 黑体谱的短波端为何受压低？以光子能量和热尺度解释，不必使用书中图片。
(c) 星际氰基 CN 分子吸收线与宇宙微波背景测温有何关系？
(d) Penzias 与 Wilson 的发现是什么，当时供职何处？
(e) Ryden 的熊与火箭示意图要表达广义相对论的哪项基本原理？说清比较的两种局部实验处境。
\end{exercise}
\begin{solution}
(a) 是均匀性与 Hubble 定律相容。若各星系在同一时刻的位置为 $\boldsymbol x_i$，流速为 $\boldsymbol v_i=H\boldsymbol x_i$，任意改以 B 为原点，仍有 $\boldsymbol v_i-\boldsymbol v_B=H(\boldsymbol x_i-\boldsymbol x_B)$；C 也一样。图中的中心只是坐标选择。
(b) 当 $\lambda$ 很短，$h\nu=hc/\lambda\gg\kb T$，占据高能光子态所需热能难以取得，Planck 因子趋于 $\exp[-hc/(\lambda\kb T)]$。它与长波端模数、体积的说法是不同的抑制机制。
(c) CN 的基态与激发态吸收线给出相对布居；在辐射与气体的适用平衡条件下，布居比约为 $(g_1/g_0)\exp[-(E_1-E_0)/(\kb T_{\rm exc})]$。由线比反推激发温度，可检验使其受激的微波背景温度；碰撞与局部辐射需单独评估，不能把任一线比自动等同 CMB 温度。
(d) Arno Penzias、Robert Wilson 发现宇宙微波背景；任职贝尔电话实验室。
(e) 是等效原理：足够小的实验室里，熊在地球重力场中静止，和在远离引力源而作相应加速的火箭中，局部力学表现可相同。有限空间或长时间测得的潮汐效应可区分，原理不宣称全球完全相同。
\end{solution}

\src{8.286 Fall 2013 Quiz 2, Q1(a)--(d), 25分；官方解答逐项可得。}
\begin{exercise}
(a) 若星系质量几乎全由中央可见恒星提供，外侧圆轨道速度 $v(R)$ 应如何随半径变化？(b) 实测典型星系外侧旋转曲线大致怎样？(c) 写出稳态、引力束缚系统的 virial 定理。(d) 解释为何早期夸克与反夸克的极小数目不对称，能够与今天显著的重子、反重子不对称及每个重子对应大量光子并存。
\end{exercise}
\begin{solution}
(a) 外侧包围质量近似常数 $M$，由 $v^2/R=GM/R^2$ 得 $v\propto R^{-1/2}$；这同时假定可近似用球对称或远场点质量。(b) 许多观测在相当长的外侧范围近乎平坦，即 $v$ 下降远慢于 $R^{-1/2}$，提示可见星体以外的质量分布；若要从速度求精确质量轮廓还需几何模型。(c) 令 $I=\sum_i m_i r_i^2$，则 $\tfrac12\ddot I=2K+W$，稳态的时间平均满足 $2\langle K\rangle+\langle W\rangle=0$；自引力 $1/r$ 势的 $W$ 等于势能。(d) 高温时夸克与反夸克数目巨大且只差一个小余数；冷却后成对湮灭，余下净夸克结合为重子，反重子极少。湮灭产物增加相对于残余重子的辐射数目。官方示意口径约每 $10^9$ 个反夸克多三个夸克，并非本题要精确计算的重子生成模型。
\end{solution}

\src{8.286 Fall 2013 Quiz 2, Q3(a)--(d), 35分；同届 Quiz 2 Review Problem 17 后重复入卷，题目在此给完整解。}
\begin{exercise}
旋转柱坐标 $(t,r,\phi,z)$ 满足 $-\infty<t,z<\infty,\ r\ge0,\ \phi\sim\phi+2\pi$，度规为
\[
\dd s^2=-c^2\dd\tau^2=-c^2\dd t^2+\dd r^2+r^2(\dd\phi+\omega\dd t)^2+\dd z^2.
\]
它来自惯性柱坐标 $\bar t=t,\bar r=r,\bar\phi=\phi+\omega t,\bar z=z$，$\omega$ 为常数。
(a) 求 $g_{tt},g_{t\phi}=g_{\phi t},g_{rr},g_{\phi\phi},g_{zz}$。
(b) 对固有时参数的测地线，利用
$\frac{\dd}{\dd\tau}(g_{\mu\nu}x'^\nu)=\frac12(\partial_\mu g_{\lambda\sigma})x'^\lambda x'^\sigma$
写 $\mu=r$ 方程。
(c) 写 $\mu=\phi$ 方程。
(d) 以坐标速度 $\dd r/\dd t,\dd\phi/\dd t,\dd z/\dd t$ 及 $r,c,\omega$ 表示未来向类时轨道的 $\dd t/\dd\tau$。
\end{exercise}
\begin{solution}
展开平方得
\[
g_{tt}=-c^2+\omega^2r^2,\quad
g_{t\phi}=\omega r^2,\quad
g_{rr}=1,\quad g_{\phi\phi}=r^2,\quad g_{zz}=1 .
\]
此处交叉项在 $\dd s^2$ 中是 $2g_{t\phi}\dd t\dd\phi$，不能误取 $2\omega r^2$ 为矩阵元。
以撇号表 $\dd/\dd\tau$，因为 $g_{r\nu}x'^\nu=r'$，
\[
r''=\tfrac12(2\omega^2r\,t'^2+4\omega r\,t'\phi'+2r\,\phi'^2)
=r(\phi'+\omega t')^2 .
\]
度规不含 $\phi$，对应共轭动量守恒，故
\[
\frac{\dd}{\dd\tau}\{r^2(\phi'+\omega t')\}=0.
\]
最后将 $\dd s^2=-c^2\dd\tau^2$ 除以 $\dd t^2$，对 $t'>0$ 得
\[
\frac{\dd t}{\dd\tau}
=\left[1-\frac{(\dd r/\dd t)^2+r^2(\dd\phi/\dd t+\omega)^2+(\dd z/\dd t)^2}{c^2}\right]^{-1/2}.
\]
根号内须正。$r=c/|\omega|$ 只是固定旋转坐标点从类时变成零向的界线，不是平坦时空的曲率奇点。官方试卷备用占位函数列表的交叉分量有笔误；直接展开给出的度规避免沿用该误印。
\end{solution}

\src{8.286 Fall 2013 Quiz 3, Q1(a)--(g), 35分；官方解答逐项可得。}
\begin{exercise}
(a) COBE/FIRAS 的黑体谱精度写作 $\Delta\epsilon/\epsilon\sim10^{-n}$，给 $n$，允许误差一位指数。
(b) 扣除前的 CMB 偶极主要为何？
(c) 扣除偶极后，以平均温度 $2.725\,\mathrm K$ 为基准，$\langle(\delta T/T)^2\rangle^{1/2}$ 约是多少，数量级即可？
(d) 以下六种解释何者对应题中 Sachs--Wolfe 效应：朝向我们运动的发光流体因 Doppler 而变红或变蓝；末次散射面过密区的光子爬出势阱而变红或变蓝；星系际介质吸收而变红或变蓝？逐项辨明物理来源。
(e) 若今天 $\Omega$ 距 1 小于约 $1\%$，传统热大爆炸模型在 $t=1\,\mathrm s$ 要求 $|\Omega-1|<10^{-m}$；给 $m$，允许误差三位。
(f) 给出题目所指 \emph{Planck 2013} 口径的今天重子、暗物质、暗能量百分比。
(g) 不含暴胀的传统模型里，末次散射面上一个粒子视界长度在天上的张角约多少度？允许因尺度定义有约两倍误差。
\end{exercise}
\begin{solution}
(a) 官方给 $n=4$；这是题目的历史灵敏度量级，不等同所有谱畸变参数的现时约束。(b) 观测者相对近似各向同性的 CMB 静止系运动造成运动学偶极，而非把它直接解释为原初密度涨落。(c) 约 $1.1\times10^{-5}$。
(d) 按原题选择是“过密区的光子向外爬出势阱而变红”；朝向我们运动的流体给蓝移，是 Doppler 项；吸收改变强度和谱形，不是本题的引力红移。完整大角标量 Sachs--Wolfe 温度扰动还含发射面本征温度与引力项，不能将选择题口头描述当作完整微扰公式。
(e) 官方量级 $m=18$；它随早期物质和辐射转换、今天允许偏离值的粗略选取而变化，关键是早期曲率偏差极小。
(f) 依 2013 年题解分别约 $5\%,26.5\%,68.5\%$，三者和为 $100\%$；这是历史题，不以今天重新拟合的参数替换原答案。
(g) 约 $2^\circ$，官方接受约 $1^\circ$--$5^\circ$。物质主导近似下粒子视界约为当时 Hubble 长度的两倍；声学视界与粒子视界不同，不能把首声学峰角尺度机械用作本题答案。
\end{solution}
